\section{From a counterexample to strict concavity}\label{sec:reduction}

\begin{idea}[What this section does]
We assume \cref{thm:main} fails and produce an oriented graph $G$ with
positive indegrees in which every nonempty proper vertex set $U$ satisfies
$\abs U+\abs{F^2U}<2\abs{FU}$ (\cref{prop:reduction}). This takes two steps.
First, a minimal counterexample satisfies a ``subset'' version of the
conjecture's failure (\cref{lem:subset-deficit}). Second, a blow-up turns
that subset version into the clean concavity inequality
(\cref{lem:blowup}). Only elementary counting is used.
\end{idea}

Throughout this section $F=F_D$ is the forward image in a fixed digraph $D$
(defined in \cref{sec:glance}); it is \emph{not} required that $F(S)$ be
disjoint from $S$.

\subsection{Minimal counterexamples}

Call an oriented graph $D$ a \emph{counterexample} if it is nonempty and
$\abs{\Nout(v)}>\abs{\Ntwo(v)}$ for every vertex $v$. A counterexample is
\emph{minimal} if it has the minimum number of vertices among all
counterexamples and, subject to this, the minimum number of arcs.

\begin{lemma}\label{lem:strong}
A minimal counterexample $D$ is strongly connected, has at least two
vertices, and every vertex has an in-neighbor.
\end{lemma}

\begin{proof}
A sink $v$ would satisfy $\abs{\Nout(v)}=0\le\abs{\Ntwo(v)}$, so $D$ has no
sink; in particular $D$ has at least two vertices. Suppose $D$ is not
strongly connected. Then some strongly connected component $K$ has no arc
leaving it (a ``sink component'' of the condensation), and $K\ne V(D)$.
For $v\in K$, all vertices reachable from $v$ lie in $K$, so $\Nout(v)$ and
$\Ntwo(v)$ are the same in $D$ and in the induced subgraph $D[K]$. Hence $D[K]$
is a counterexample with fewer vertices, which is impossible. So $D$ is
strongly connected. Since $D$ has at least two vertices, every vertex then
has an in-neighbor.
\end{proof}

\subsection{The subset deficit}\label{sec:subset-deficit}

The defining property of a counterexample concerns single vertices. The
next lemma upgrades it to all nonempty proper sets, at the price of
replacing $\Nout$ and $\Ntwo$ by the set analogues $F(S)\setminus S$ and
$F^2(S)\setminus F(S)$. For $S=\{v\}$ these are exactly $\Nout(v)$ and
$\Ntwo(v)$, so the lemma generalizes the counterexample property.

\begin{lemma}[Subset deficit]\label{lem:subset-deficit}
Let $D$ be a minimal counterexample. Then
\begin{equation}\label{eq:subset-deficit}
 \abs{F^2(S)\setminus F(S)}<\abs{F(S)\setminus S}
 \qquad\text{for every }\varnothing\ne S\subsetneq V(D).
\end{equation}
\end{lemma}

The argument adapts the proof of~\cite[Lemma~4]{Seacrest2019}, where the
right-hand side is $\abs{F(S)}$. Here we need the stronger right-hand side
$\abs{F(S)\setminus S}$, and this is why we minimize the number of vertices
before the number of arcs (to get strong connectivity).

\paragraph{Intuition.}
Fix $S$. Split $F(S)$ into the part $I=F(S)\cap S$ inside $S$ and the part
$E=F(S)\setminus S$ outside. Since $F(I)\subseteq F(S)$, everything in
$F^2(S)\setminus F(S)$ is reached from $E$. We ``switch on'' the vertices
of $E$ in batches. Each time a batch $Q$ is switched on, the set of vertices
reached from the switched-on part grows, and we show that it grows by
fewer than $\abs Q$ vertices. The batch $Q$ comes from a Seymour vertex
$v\in S$ of a graph with fewer arcs, and the growth is bounded by comparing
what $v$ has in $D$ with what it has in the smaller graph.

\begin{figure}[ht]
\centering
\begin{tikzpicture}[scale=0.95]
 % S
 \draw[thick,fill=famL!8] (0,0) ellipse (1.6 and 2.3);
 \node[famL!80!black] at (-1.1,2.25) {$S$};
 \draw[fill=famL!22,draw=famL] (0,1.05) ellipse (1.05 and 0.75);
 \node at (0,1.05) {$I$};
 \node[vtx,fill=black!10] (v) at (0,-1.2) {$v$};
 % E
 \draw[thick,rounded corners,fill=zcol!7] (3.3,-2.3) rectangle (5.7,2.3);
 \node[zcol!70!black] at (4.5,2.65) {$E=F(S)\setminus S$};
 \draw[rounded corners,fill=zcol!20] (3.5,0.85) rectangle (5.5,2.1);
 \node at (4.5,1.47) {$T$};
 \draw[rounded corners,fill=famR!25,draw=famR] (3.5,-0.6) rectangle (5.5,0.6);
 \node at (4.5,0) {$Q$};
 \node[font=\small,align=center,text=black!60] at (4.5,-1.45) {rest of\\ $E\setminus T$};
 % C
 \draw[rounded corners,fill=hcol!12,draw=hcol] (7.6,-1.2) rectangle (9.4,1.2);
 \node[align=center] at (8.5,0) {$C$};
 % arcs
 \draw[arc,red!70!black,densely dashed] (v) to[out=30,in=200] (3.55,-0.2);
 \draw[arc,red!70!black,densely dashed] (v) to[out=-5,in=190] (3.55,-1.5);
 \draw[arc] (v) to[out=60,in=230] (3.55,1.2);
 \draw[arc] (v) to[out=100,in=-90] (-0.2,0.35);
 \draw[arc] (5.5,0.1) -- (7.6,0.1);
 \draw[arc] (5.5,1.6) to[out=0,in=150] (7.0,2.0) node[right,font=\small] {$F(T)$};
 \node[font=\footnotesize,red!70!black,align=center] at (1.0,-2.95) {deleted in $D_T$: all arcs $S\to E\setminus T$};
\end{tikzpicture}
\caption{One step in the proof of \cref{lem:subset-deficit}, with
$I=F(S)\cap S$. All arcs from
$S$ to $E\setminus T$ are deleted (dashed). The minimal choice of $D$ gives a
Seymour vertex $v$ of the smaller graph $D_T$. It must lie in $S$ and lose
its out-neighbors $Q\subseteq E\setminus T$. The set $C$ consists of vertices
that $v$ reached in two steps only through $Q$. Since $v$ is not a Seymour
vertex of $D$, $\abs C-\abs{Q\cap F(T)}<\abs Q$.}
\label{fig:deletion}
\end{figure}

\begin{proof}[Proof of \cref{lem:subset-deficit}]
Fix $\varnothing\ne S\subsetneq V(D)$ and set
\[
 I=F(S)\cap S,\qquad E=F(S)\setminus S,\qquad
 g(T)=\abs{F(T)\setminus(I\cup T)}\quad\text{for }T\subseteq E.
\]
Since $D$ is strongly connected (\cref{lem:strong}) and $S$ is a nonempty
proper subset, some arc leaves $S$; hence $E\ne\varnothing$. Also
$g(\varnothing)=0$.

\emph{Step 1: the target is $g(E)$.} Since $F(S)=I\cup E$ and
$F(I)\subseteq F(S)$ (because $I\subseteq S$),
\[
 F^2(S)\setminus F(S)=\bigl(F(I)\cup F(E)\bigr)\setminus(I\cup E)
 =F(E)\setminus(I\cup E),
\]
so $\abs{F^2(S)\setminus F(S)}=g(E)$, and we must show $g(E)<\abs E$.

\emph{Step 2: one batch.} We claim that for every $T\subsetneq E$ there is a
nonempty $Q\subseteq E\setminus T$ with
\begin{equation}\label{eq:increment}
 g(T\cup Q)-g(T)<\abs Q.
\end{equation}
Let $D_T$ be obtained from $D$ by deleting all arcs from $S$ to $E\setminus T$
(\cref{fig:deletion}). Every vertex of $E\setminus T\neq\varnothing$ has an
in-neighbor in $S$, so at least one arc is deleted. $D_T$ has the same
vertices and fewer arcs than $D$, so by minimality it is not a
counterexample: some vertex $v$ satisfies
$\abs{N^+_{1,D_T}(v)}\le\abs{N^+_{2,D_T}(v)}$.

\begin{itemize}[leftmargin=*]
\item \emph{$v$ lost an out-neighbor.} Otherwise $N^+_{1,D_T}(v)=\Nout(v)$, and
since deleting arcs can only shrink $F^2(\{v\})$, also
$N^+_{2,D_T}(v)\subseteq\Ntwo(v)$. Then $v$ would be a Seymour vertex of $D$.
Hence $v\in S$ and
\[
 Q:=\Nout(v)\cap(E\setminus T)\ne\varnothing
\]
is exactly the set of out-neighbors that $v$ lost. As
$\Nout(v)\subseteq F(S)=I\cup E$, we get $\Nout(v)\subseteq I\cup T\cup Q$
and $N^+_{1,D_T}(v)=\Nout(v)\setminus Q\subseteq I\cup T$.

\item \emph{Where $v$ can go in two steps in $D_T$.} A path $v\to u\to w$ in
$D_T$ has $u\in I\cup T$. If $u\in I\subseteq S$, then $w\in F(S)$ but $w\notin
E\setminus T$ (those arcs were deleted), so $w\in I\cup T$. If $u\in T$,
then $w\in F(T)$. Thus
\begin{equation}\label{eq:DT-reach}
 F^2_{D_T}(\{v\})\subseteq I\cup T\cup F(T).
\end{equation}

\item \emph{Gains.} A vertex $w\in N^+_{2,D_T}(v)\setminus\Ntwo(v)$ is reached
in two steps in $D_T$, hence in $D$, and $w\ne v$; so the only way it was
not in $\Ntwo(v)$ is that $w\in\Nout(v)$. As $w\notin N^+_{1,D_T}(v)$, we get
$w\in Q$, and by~\eqref{eq:DT-reach} and $Q\cap(I\cup T)=\varnothing$,
$w\in Q\cap F(T)$. So at most $\abs{Q\cap F(T)}$ vertices are gained.

\item \emph{Losses.} Put $C=F(Q)\setminus(I\cup T\cup Q\cup F(T))$. Every
$c\in C$ is reached from $v$ via some $q\in Q$. It is not an out-neighbor of
$v$, because $\Nout(v)\subseteq I\cup T\cup Q$. And $c\ne v$, since an arc
$q\to v$ would be opposite to $v\to q$ (\emph{here orientation is used}).
So $C\subseteq\Ntwo(v)$. By~\eqref{eq:DT-reach}, no vertex of $C$ is in
$F^2_{D_T}(\{v\})$. So all of $C$ is lost.
\end{itemize}
Writing $d_1=\abs{\Nout(v)}$ and $d_2=\abs{\Ntwo(v)}$, the Seymour property of
$v$ in $D_T$ and the gains/losses give
\[
 d_1-\abs Q=\abs{N^+_{1,D_T}(v)}\le\abs{N^+_{2,D_T}(v)}\le d_2-\abs C+\abs{Q\cap F(T)}.
\]
On the other hand, the increment of $g$ is exactly
\begin{equation}\label{eq:g-increment}
 g(T\cup Q)-g(T)=\abs C-\abs{Q\cap F(T)}.
\end{equation}
Indeed $F(T\cup Q)\setminus(I\cup T\cup Q)$ is the disjoint union of
$F(T)\setminus(I\cup T\cup Q)$ and $C$, and the first set is
$F(T)\setminus(I\cup T)$ with the points of $Q\cap F(T)$ removed (recall
$Q\cap(I\cup T)=\varnothing$). Combining, and using $d_2<d_1$ because $D$ is a
counterexample,
\[
 g(T\cup Q)-g(T)=\abs C-\abs{Q\cap F(T)}\le\abs Q+d_2-d_1<\abs Q,
\]
which is~\eqref{eq:increment}.

\emph{Step 3: summing.} Start from $T_0=\varnothing$ and, as long as
$T_k\ne E$, apply Step~2 (always to the \emph{original} $D$) to get
$Q_k$ and set $T_{k+1}=T_k\cup Q_k$. The sets $Q_k$ are nonempty and
disjoint, so the process ends with $T_K=E$ after $K\ge1$ steps (as
$E\neq\varnothing$). Summing~\eqref{eq:increment},
\[
 g(E)=g(\varnothing)+\sum_{k<K}\bigl(g(T_{k+1})-g(T_k)\bigr)<\sum_{k<K}\abs{Q_k}=\abs E.
\]
By Step~1 this is~\eqref{eq:subset-deficit}.
\end{proof}

\subsection{Why the deficit is not yet enough}

The subset deficit has a nice form, but we will need to \emph{add} the
inequality over many sets, and for that the following cleaner quantity is
better suited:
\[
 \kappa(S)=\abs S-2\abs{F(S)}+\abs{F^2(S)}
 =\bigl(\abs{F^2S}-\abs{FS}\bigr)-\bigl(\abs{FS}-\abs S\bigr),
\]
the ``second difference'' of the sequence $\abs S,\abs{FS},\abs{F^2S}$. It is
a plain combination of cardinalities with no set differences, so it adds
up well (we will use this in \cref{eq:fiber-sum}).

\begin{remark}[The error term]\label{rem:error-term}
For every digraph and every set $S$,
\begin{equation}\label{eq:error-term}
 \kappa(S)\le\bigl(\abs{F^2S\setminus FS}-\abs{FS\setminus S}\bigr)+\abs{S\setminus FS}.
\end{equation}
Indeed, $\abs{F^2S}-\abs{FS}\le\abs{F^2S\setminus FS}$ and
$\abs{FS}-\abs S=\abs{FS\setminus S}-\abs{S\setminus FS}$. So the subset deficit
gives $\kappa(S)\le -1+\abs{S\setminus F(S)}$, which is useless when many
vertices of $S$ receive no arc from $S$. The blow-up below multiplies the
bracket in~\eqref{eq:error-term} by a large factor $m$ while increasing the
error term by at most one per vertex of $D$.
\end{remark}

\subsection{Blowing up by transitive tournaments}\label{sec:blowup}

Let $T_m$ be the \emph{transitive tournament} on $\{1,\dots,m\}$, with
$a\to b$ if and only if $a<b$. We write $f$ for its forward image.

\begin{lemma}[Internal estimate]\label{lem:internal}
For every nonempty $T\subseteq\{1,\dots,m\}$,
\[
 \abs T-2\abs{f(T)}+\abs{f^2(T)}\le 1 .
\]
\end{lemma}

\begin{proof}
Let $a=\min T$. Then $f(T)=\{a+1,\dots,m\}$ and $f^2(T)=\{a+2,\dots,m\}$,
while $\abs T\le m-a+1$. If $a<m$ the left side is at most
$(m-a+1)-2(m-a)+(m-a-1)=0$. If $a=m$ then $T=\{m\}$ and the left side is
$1-0+0=1$.
\end{proof}

For example, in $T_5$ the set $T=\{2,4\}$ has $f(T)=\{3,4,5\}$ and
$f^2(T)=\{4,5\}$, giving $2-6+2=-2\le1$. The point of the transitive
tournament is that a nonempty set loses at most one element under $f$
($\abs{f(T)}\ge\abs T-1$). If we blew up by independent sets instead, a block
would have empty image and the error would be as large as $m$.

\begin{definition}[Blow-up]
Let $D$ be an oriented graph and $m\ge1$. The \emph{blow-up} $G=D[T_m]$ has
vertex set $V(D)\times\{1,\dots,m\}$. Its arcs are
\[
 (v,a)\to(w,b)\iff v\to w\text{ in }D,\ \text{or}\ v=w\text{ and }a<b.
\]
The set $\{v\}\times\{1,\dots,m\}$ is the \emph{block} of $v$. Arcs inside a
block are \emph{internal} and arcs between blocks are \emph{external}.
\end{definition}

\begin{figure}[ht]
\centering
\begin{tikzpicture}[scale=0.95]
 \begin{scope}
  \node[vtx] (a) at (90:1.2) {$x$};
  \node[vtx] (b) at (210:1.2) {$y$};
  \node[vtx] (c) at (330:1.2) {$z$};
  \draw[arc] (a)--(c); \draw[arc] (c)--(b); \draw[arc] (b)--(a);
  \node at (0,-2.0) {$D$};
 \end{scope}
 \begin{scope}[xshift=6.0cm]
  \foreach \nm/\ang in {x/90,y/210,z/330}{
   \begin{scope}[shift={(\ang:2.0)}]
    \draw[rounded corners,fill=black!4] (-0.45,-1.0) rectangle (0.45,1.0);
    \node[dot] (\nm1) at (0,0.65) {}; \node[dot] (\nm2) at (0,0) {}; \node[dot] (\nm3) at (0,-0.65) {};
    \draw[-{Stealth[length=1.5mm]}] (\nm1)--(\nm2);
    \draw[-{Stealth[length=1.5mm]}] (\nm2)--(\nm3);
    \draw[-{Stealth[length=1.5mm]}] (\nm1) to[bend right=55] (\nm3);
    \node[font=\small] at (0,1.25) {$\nm$};
   \end{scope}}
  \draw[-{Stealth[length=3mm]},line width=2.2pt,black!35] (90:2.0)++(0.55,-0.4) -- ($(330:2.0)+(-0.25,0.95)$);
  \draw[-{Stealth[length=3mm]},line width=2.2pt,black!35] ($(330:2.0)+(-0.55,-0.6)$) -- ($(210:2.0)+(0.55,-0.6)$);
  \draw[-{Stealth[length=3mm]},line width=2.2pt,black!35] ($(210:2.0)+(0.25,0.95)$) -- ($(90:2.0)+(-0.55,-0.4)$);
  \node at (0,-2.75) {$G=D[T_3]$};
 \end{scope}
 \begin{scope}[xshift=11.3cm,yshift=1.2cm]
  \foreach \k in {1,...,5} \node[vtx,minimum size=5mm,font=\footnotesize] (t\k) at (0,-0.75*\k+0.75) {\k};
  \foreach \k in {2,4} \node[vtx,minimum size=5mm,font=\footnotesize,fill=famL!35] at (t\k) {\k};
  \node[font=\small,align=left,anchor=west] at (0.5,-0.4) {$T=\{2,4\}$};
  \node[font=\small,align=left,anchor=west] at (0.5,-1.2) {$f(T)=\{3,4,5\}$};
  \node[font=\small,align=left,anchor=west] at (0.5,-2.0) {$f^2(T)=\{4,5\}$};
  \node at (0,-3.95) {in $T_5$};
 \end{scope}
\end{tikzpicture}
\caption{Left: a base graph $D$ (a directed triangle). Middle: its blow-up
$D[T_3]$. Each vertex becomes a transitive tournament, and each arc becomes
all $9$ arcs between the corresponding blocks (thick grey arrows). Right:
images inside a transitive tournament (vertices ordered top to bottom; every
arc points downwards).}
\label{fig:blowup}
\end{figure}

\begin{lemma}[Blow-up]\label{lem:blowup}
Let $D$ be an oriented graph on $n\ge1$ vertices in which every vertex has
an in-neighbor and which satisfies~\eqref{eq:subset-deficit}. Let $m>n$ and
$G=D[T_m]$. Then $G$ is an oriented graph, every vertex of $G$ has an
in-neighbor, and
\begin{equation}\label{eq:strict-concavity}
 \abs U+\abs{F_G^2(U)}<2\abs{F_G(U)}\qquad\text{for every }\varnothing\ne U\subsetneq V(G).
\end{equation}
\end{lemma}

The same blow-up appears in~\cite[Theorem~4.3]{BrantnerBrockmanKaySnively2009},
where it is applied to single vertices. Here we need the estimate for all
sets.

\begin{proof}
$G$ is oriented: internal arcs go from smaller to larger labels, and two
opposite external arcs would come from opposite arcs of $D$. Each $(w,b)$
receives an arc from the block of an in-neighbor of $w$ in $D$.

Fix $\varnothing\ne U\subsetneq V(G)$. For $v\in V(D)$ let
$T_v=\{a:(v,a)\in U\}$ and let $S=\{v:T_v\ne\varnothing\}$ be the
\emph{support} of $U$. We compute block by block.

\emph{First image (exact).} An external arc into the block of $w$ starting
in $U$ exists only if $w\in F(S)$, and then it reaches \emph{every} vertex of
the block. Internal arcs from $U$ into the block of $w$ reach exactly
$\{w\}\times f(T_w)$. Hence the block of $w$ meets $F_G(U)$ in the whole block
if $w\in F(S)$, in $\{w\}\times f(T_w)$ if $w\in S\setminus F(S)$, and not at all
otherwise:
\begin{equation}\label{eq:first-image}
 \abs{F_G(U)}=m\abs{F(S)}+\sum_{w\in S\setminus F(S)}\abs{f(T_w)}.
\end{equation}

\emph{Second image (upper bound).} Consider a path of length two starting in
$U$. If both arcs are external, it ends in a block of $F^2(S)$. If exactly
one is external, it ends in a block of $F(S)$. If the internal step comes
first, the path stays in a block of $S$ and then moves to a block of $F(S)$.
If the external step comes first, it moves to a block of $F(S)$ and the
internal step stays in that block. If
both arcs are internal, it stays in the block of some $v\in S$ and ends in
$\{v\}\times f^2(T_v)$. Writing $H_0=F(S)\cup F^2(S)$, and counting the blocks
of $H_0$ in full, we get
\begin{equation}\label{eq:second-image}
 \abs{F_G^2(U)}\le m\abs{H_0}+\sum_{v\in S\setminus H_0}\abs{f^2(T_v)}
 \le m\abs{H_0}+\sum_{v\in S\setminus F(S)}\abs{f^2(T_v)},
\end{equation}
where the last step only adds nonnegative terms (as $F(S)\subseteq H_0$).

\emph{Combining.} Split $\abs U=\sum_{v\in S\cap F(S)}\abs{T_v}+\sum_{v\in S\setminus F(S)}\abs{T_v}$.
From~\eqref{eq:first-image} and~\eqref{eq:second-image},
\[
 \abs U-2\abs{F_GU}+\abs{F_G^2U}
 \le m\abs{H_0}-2m\abs{F(S)}+\sum_{v\in S\cap F(S)}\abs{T_v}
 +\sum_{v\in S\setminus F(S)}\Bigl(\abs{T_v}-2\abs{f(T_v)}+\abs{f^2(T_v)}\Bigr).
\]
Use $\abs{T_v}\le m$ in the first sum and \cref{lem:internal} in the second:
\[
 \abs U-2\abs{F_GU}+\abs{F_G^2U}\le m\bigl(\abs{H_0}-2\abs{F(S)}+\abs{S\cap F(S)}\bigr)+\abs{S\setminus F(S)}.
\]
Since $\abs{H_0}=\abs{F(S)}+\abs{F^2(S)\setminus F(S)}$ and
$\abs{F(S)}-\abs{S\cap F(S)}=\abs{F(S)\setminus S}$, the bracket equals
$\abs{F^2(S)\setminus F(S)}-\abs{F(S)\setminus S}$. We have shown, for every
$U$,
\begin{equation}\label{eq:blowup-master}
 \abs U-2\abs{F_GU}+\abs{F_G^2U}\le m\bigl(\abs{F^2S\setminus FS}-\abs{FS\setminus S}\bigr)+\abs{S\setminus FS},
\end{equation}
which is the scaled version of~\eqref{eq:error-term} announced in
\cref{rem:error-term}.

\emph{Case $S\ne V(D)$.} By~\eqref{eq:subset-deficit} the bracket is an integer
$\le-1$, and $\abs{S\setminus F(S)}\le n$, so the right side is at most
$-m+n<0$.

\emph{Case $S=V(D)$.} Every vertex of $D$ has an in-neighbor, so
$F(S)=V(D)$ and every block is filled: $F_G(U)=V(G)$. Since every vertex of
$G$ has an in-neighbor, $F_G^2(U)=F_G(V(G))=V(G)$ too. Thus the left side
equals $\abs U-\abs{V(G)}<0$, as $U$ is proper.
\end{proof}

\begin{proposition}[Counterexample reduction]\label{prop:reduction}
If \cref{thm:main} fails, there is a nonempty oriented graph $G$ in which
every vertex has an in-neighbor and~\eqref{eq:strict-concavity} holds for
every nonempty proper vertex set $U$.
\end{proposition}

\begin{proof}
Take a minimal counterexample $D$. By \cref{lem:strong,lem:subset-deficit}
it has positive indegrees and satisfies~\eqref{eq:subset-deficit}. Apply
\cref{lem:blowup} with any $m>\abs{V(D)}$.
\end{proof}

\begin{idea}[What we have achieved, and what is used later]
From now on we only use the graph $G$ and two of its properties:
\textbf{(G1)} $G$ is a nonempty oriented graph in which every vertex has an
in-neighbor; \textbf{(G2)} every nonempty proper $U\subseteq V(G)$ is
strictly concave, $\abs U+\abs{F^2U}<2\abs{FU}$. The remaining sections show
that (G1) and (G2) are incompatible.
\end{idea}
